\subsection{Study B4 — Grouped Lexical Residualization}
\label{sec:results-study-b4}

Study B3's cross-validated residualization performed random $70$/$30$
splits of the 600-probe set. Random splits are honest at the prompt-row
level but leak template-family structure: held-out rows almost always
share a template family with training rows, because each template was
sampled approximately uniformly across the random split. The review
of the previous version of this manuscript correctly flagged this as
the next decisive control to run.

Study B4 is the grouped cross-validation. Three grouping schemes are
implemented:

\begin{itemize}
    \item \textbf{Leave-One-Template-Family-Out (LOTFO):} the 600
        probes are partitioned into $28$ template families
        (recovered post-hoc by exact substring matching against the
        deterministic generators in
        \texttt{scripts/build\_phase1\_probes.py}). For each
        template family, fit ridge on the other $27$ families,
        residualize the held-out family's prompts, compute metrics
        on those residuals.
    \item \textbf{Leave-One-Domain-Out (LODO):} hold out one of
        \{math, logic, polecon\} entirely; train on the other two.
    \item \textbf{Two-vs-one split:} train on math$+$logic, test on
        polecon (the largest single grouped split available).
\end{itemize}

For each grouped split, training-time TF-IDF and template-only
vocabularies are fit on the training rows only; held-out test rows
are featurized using the training vocabulary and IDF weights, with
out-of-vocabulary test tokens dropped. This is the proper way to
extend lexical featurization to grouped held-out rows.

\paragraph{Result.} Table~\ref{tab:study-b4-lotfo} reports LOTFO.
Table~\ref{tab:study-b4-lodo} reports LODO and the two-vs-one split.

\begin{table}[h]
\centering
\caption{Study B4 LOTFO ($28$ leave-one-template-family-out splits).
Variance-removed $95\%$ percentile intervals frequently span zero
and reach substantially negative values --- the lexical regression
fitted on $27$ template families does not generalize to the
held-out family. CKA is essentially unchanged by residualization:
$\Delta_\text{CKA} = -0.006, +0.019, +0.010$ across layers.}
\label{tab:study-b4-lotfo}
\begin{tabular}{l c c c c c}
\toprule
Layer & Var.\ removed (A) & Var.\ removed (B) & CKA orig & CKA resid & ΔCKA \\
\midrule
$0.25$ & $0.32$ [$-0.44$, $0.71$] & $0.22$ [$-0.50$, $0.59$] & $0.808$ & $0.802$ & $-0.006$ \\
$0.50$ & $0.00$ [$-0.90$, $0.54$] & $0.06$ [$-0.78$, $0.53$] & $0.705$ & $0.723$ & $+0.019$ \\
$0.75$ & $-0.01$ [$-0.98$, $0.59$] & $0.01$ [$-1.07$, $0.58$] & $0.600$ & $0.610$ & $+0.010$ \\
\bottomrule
\end{tabular}
\end{table}

\begin{table}[h]
\centering
\caption{Study B4 LODO and two-vs-one split. Held-out variance
removed is uniformly negative (the lexical regression makes
worse-than-mean predictions on the held-out group), and CKA
\emph{increases} after residualization in every cell, by as much as
$+0.34$. The lexical residualization fails outright when the held-out
prompts share no template family with the training prompts.}
\label{tab:study-b4-lodo}
\begin{tabular}{l l c c c c c c}
\toprule
Held-out & Layer & Var.\ rem.\ (A) & Var.\ rem.\ (B) & CKA orig & CKA resid & ΔCKA & ΔProc \\
\midrule
\multicolumn{8}{l}{\emph{Leave-One-Domain-Out}} \\
logic   & $0.25$ & $-0.49$ & $-0.45$ & $0.816$ & $0.858$ & $+0.042$ & $-0.006$ \\
math    & $0.25$ & $-1.12$ & $-1.20$ & $0.904$ & $0.948$ & $+0.043$ & $-0.024$ \\
polecon & $0.25$ & $-0.31$ & $-0.61$ & $0.774$ & $0.815$ & $+0.041$ & $-0.033$ \\
logic   & $0.50$ & $-0.64$ & $-0.46$ & $0.668$ & $0.770$ & $+0.103$ & $-0.002$ \\
math    & $0.50$ & $-1.93$ & $-1.29$ & $0.797$ & $0.914$ & $+0.117$ & $+0.005$ \\
polecon & $0.50$ & $-1.41$ & $-1.06$ & $0.556$ & $0.891$ & $+0.335$ & $-0.015$ \\
logic   & $0.75$ & $-0.78$ & $-0.44$ & $0.480$ & $0.621$ & $+0.141$ & $-0.001$ \\
math    & $0.75$ & $-1.82$ & $-1.44$ & $0.635$ & $0.820$ & $+0.186$ & $+0.001$ \\
polecon & $0.75$ & $-1.37$ & $-0.98$ & $0.558$ & $0.859$ & $+0.301$ & $-0.011$ \\
\midrule
\multicolumn{8}{l}{\emph{Two-vs-one (train math+logic, test polecon)}} \\
polecon & $0.25$ & $-0.31$ & $-0.61$ & $0.774$ & $0.815$ & $+0.042$ & $-0.033$ \\
polecon & $0.50$ & $-1.41$ & $-1.06$ & $0.556$ & $0.891$ & $+0.335$ & $-0.015$ \\
polecon & $0.75$ & $-1.37$ & $-0.98$ & $0.558$ & $0.859$ & $+0.301$ & $-0.011$ \\
\bottomrule
\end{tabular}
\end{table}

\paragraph{What Study B4 establishes.} The grouped cross-validation
forces a substantial reframing of the negative methodological finding
from Studies B2, B3.

\begin{enumerate}
    \item \textbf{The lexical structure that random-split
        residualization removed is template-specific.} The TF-IDF
        and template-only featurization fit on $27$ template
        families does not generalize to the held-out template family:
        variance-removed point estimates hover near zero with $95\%$
        percentile intervals that extend substantially negative.
        Under leave-one-domain-out the regression is even worse
        (variance-removed strongly negative across the board: the
        lexical model fitted on two domains makes worse-than-mean
        predictions on the held-out domain).
    \item \textbf{Under grouped CV, lexical residualization barely
        affects CKA --- and often \emph{increases} it.} LOTFO
        $\Delta_\text{CKA}$ ranges from $-0.006$ to $+0.019$
        across layers (essentially zero). LODO and two-vs-one
        $\Delta_\text{CKA}$ are positive in every cell, ranging from
        $+0.04$ to $+0.34$. The CKA increase is partly a consequence
        of the residualization failure: when ridge regression
        produces worse-than-mean predictions, the residuals contain
        correlated noise that, when present in both models, can
        \emph{increase} apparent similarity. We therefore do not read
        the LODO/two-vs-one $\Delta_\text{CKA} > 0$ as evidence of
        ``lexical residualization strengthening alignment''; we read
        it as ``lexical residualization fails outright in this
        regime.''
    \item \textbf{Procrustes residual is essentially unchanged
        under grouped CV.} The Procrustes residual changes are at
        most $\pm 0.04$ in any cell of LOTFO/LODO/two-vs-one. Combined
        with the residualization failure, this is consistent with
        ``the lexical regression is too template-specific to remove
        meaningful variance from out-of-template activations.''
    \item \textbf{The lexical contribution to cross-model alignment is
        not as large as Studies B2/B3 suggested.} The B3
        cross-validated residualization removed a substantial
        portion of CKA ($-0.37$ to $-0.49$ on held-out random splits).
        Under grouped CV the same residualization removes essentially
        none of CKA. The honest reading is that B3's effect was
        partly an artifact of within-template lexical structure that
        random splits leak across train and test.
\end{enumerate}

\paragraph{Methodological lesson, expanded.} Random-split lexical
controls in cross-model representation studies overstate the
contribution of lexical structure to apparent similarity. Grouped
cross-validation is the necessary strict version. But Study B4 also
shows the limits of TF-IDF-based lexical featurization itself: a
$455$-dimensional bag-of-words featurization does not generalize
across template families in a way that supports clean residualization
on held-out templates. Future cross-model representation studies that
use lexical residualization as a control should: (i) use
domain-general lexical features (subword tokenizer counts,
pre-trained sentence-encoder representations, character-n-gram TF-IDF
fitted on a large corpus rather than the probe set itself),
(ii) report grouped cross-validation by template family and domain
explicitly, (iii) report the variance-removed statistic for the
held-out group to verify the lexical regression actually generalizes
before interpreting the residualization result.

\paragraph{What this means for the paper's central claim.} The
strong reading of Studies B1/B2/B3 was: most of the cross-model
representational similarity at these layers is shared lexical and
template encoding. Study B4 weakens this claim. The honest reading
after Study B4 is: \emph{a substantial portion of the cross-model
similarity is template-specific encoding that random-split
residualization can remove, but the structure that survives a
template-grouped lexical control is much larger than the B3 number
indicated.} The cross-model alignment is more robust to grouped
lexical control than B3 suggested. This does not rescue the strong
operator-theoretic interpretation: the surviving signal could still
be training-distribution residue, deeper syntactic structure, or
other shared encoding that our lexical featurization fails to capture.
But the negative finding from B1/B2/B3 must be read with the B4
qualification: random-split lexical controls overstate the lexical
contribution.
